I am applying for the position of \textit{Postdoctoral Fellow in Cyber Range Applications for Cyber Security} within the Critical Infrastructure Security and Resilience (CISaR) group at NTNU. I am completing my PhD in Cybersecurity at UiT The Arctic University of Norway, where my research focuses on device-level protection through lightweight cryptography and system-level trust enforcement via blockchain-based access control. This position aligns closely with my doctoral work on secure system design and provides an opportunity to extend that expertise toward large-scale testing, validation, and cyber range experimentation for resilient infrastructures.

I am particularly motivated by NTNU’s leadership in critical infrastructure protection and applied cybersecurity research within the CISaR group. I have followed your recent publication, \textit{“A Step-by-Step Definition of a Reference Architecture for Cyber Ranges”} (Computers \& Security, 2024), which establishes a structured methodology for designing and evaluating cyber range environments. The paper’s emphasis on architectural modularity, interoperability, and automated validation resonates deeply with my research notion of combining cryptographic assurance with adaptive system validation. In addition, NTNU’s involvement in the Horizon Europe projects \textit{SECASSURED} and \textit{ATHENA}, which aim to functionalize advanced validation frameworks across critical infrastructure sectors, presents an ideal environment to apply my background in trustworthy system design to practical, testbed-based research.

During my PhD, I developed the Distributed Trust Enforcement System (DTES), a blockchain-based framework for dynamic policy enforcement and capability-based access control in distributed IoT systems. My work involved security modeling, protocol design, and empirical evaluation using lightweight cryptographic algorithms on constrained hardware. Building on this experience, I am eager to expand my expertise into cyber range–driven validation and AI-assisted security testing. I bring strong programming skills in C/C++, Python, and Solidity, familiarity with automated validation pipelines, and a collaborative mindset suited for interdisciplinary research within CISaR’s projects and partnerships.

Thank you for considering my application.